\section*{Scriptural Date and Location}

\begin{dossiertable}
\dossierrow{Entrance Antiphon}{Psalm 54:6, 8 \newline (Vulgate 53:6, 8)}{Judahite temple prayer; transmitted in the Psalter of Jerusalem worship}{Final form uncertain; collection closed in the Second Temple period}
\dossierevent{1 Samuel 23:19; 26:1}{Wilderness of Ziph, Judean hill country south-east of Hebron}{Reign of Saul, traditionally eleventh century BC}
\dossiernote{The superscription assigns the psalm to David when the Ziphites betrayed his hiding place; inherited attribution is not a dated autograph, and historical judgment leaves author and date open. The two half-verses the Missal prints stand at the psalm's turn from danger to willing sacrifice. The Missal cites the psalm as 53 in the Vulgate series; the same text is numbered 54 in the Hebrew series used by the United States Lectionary.}

\dossierrow{Responsorial Psalm}{Psalm 86:5--6, 9--10, 15--16; response from 86:5a}{An individual's plea, composed largely of phrases drawn from older psalms and from Exodus 34:6}{Uncertain; the borrowing suggests a relatively late compilation}
\dossiernote{The psalm carries a Davidic heading but no narrated occasion. Verse 9---all the nations coming to worship---places an individual petition inside a universal horizon; verse 15 quotes the divine self-description of Exodus 34:6. The appointed verses omit the psalm's complaint against violent adversaries (86:14).}

\dossierrow{Communion Antiphon, first}{Psalm 111:4--5 \newline (Vulgate 110:4--5)}{Praise sung in the assembly of the upright (111:1); temple setting}{Post-exilic; commonly placed in the Persian or early Hellenistic period}
\dossiernote{An alphabetic acrostic in Hebrew: each half-line begins with a successive letter, a form that survives in neither Latin nor English. The remembered wonders are the Exodus deeds; the food given is, in the psalm's own frame, the covenant provision, not yet the Eucharist.}

\dossierrow{First Reading}{Wisdom 12:13, 16--19}{Written in Greek at Alexandria in Egypt for a Jewish community living under Hellenistic and then Roman rule}{Commonly placed in the last century BC}
\dossierevent{Wisdom 12:3--11}{The land of Canaan, in the book's retrospect on its earlier inhabitants}{Israel's occupation of the land, in the book's own frame}
\dossiernote{The book speaks in the persona of Solomon and is preserved only in Greek; Catholic canons receive it, and it is not in the Hebrew Bible. The appointed verses belong to a defence of God's slowness against Canaan; the Lectionary omits verses 14--15, which deny that any king or tyrant may call God to account and state that God's justice is the reason he governs justly.}

\dossierrow{Gospel Acclamation}{Cf.\ Matthew 11:25}{Adapted, not quoted: Jesus' thanksgiving in Galilee, before the parable discourse}{See the Gospel dossier below}
\dossiernote{The verse is a liturgical adaptation, marked \incipit{Cf.}\ in the Lectionary; it compresses Matthew 11:25 and supplies the phrase \emph{mysteries of the kingdom}, which in Matthew belongs to 13:11 rather than to 11:25. It is an acclamation greeting the Gospel, not a passage read in course.}

\dossierrow{Gospel (both forms)}{Matthew 13:24--43; shorter form 13:24--30}{Composed in Greek for a largely Jewish-Christian readership; Syria, and Antioch in particular, is the common modern proposal}{Ancient tradition names the apostle Matthew; the modern majority dates the Greek Gospel after AD 70, often c.\ AD 80--90}
\dossierevent{Matthew 13:1, 36}{Beside the Sea of Galilee, then indoors in a house---two distinct settings inside one appointed passage}{Galilean ministry, c.\ AD 28--30}
\dossiernote{The longer form crosses from public parable to private explanation; the shorter form stops inside the first parable and never leaves the shore. At 13:35 Matthew cites Psalm 78:2 (Vulgate 77:2) as fulfilled; that quotation is part of Matthew's argument and is not separately appointed.}

\dossierrow{Second Reading}{Romans 8:26--27}{Written from Corinth to house churches at Rome of mixed Jewish and Gentile membership}{c.\ AD 56--58}
\dossiernote{Undisputed Pauline authorship; the letter names Tertius as the scribe (16:22). These two verses continue the groaning sequence of 8:22--23 and are read in course, not chosen to comment on the Gospel.}

\dossierrow{Communion Antiphon, second}{Revelation 3:20}{Written on Patmos; addressed to the church at Laodicea in the Lycus valley of the Roman province of Asia}{Traditionally the last years of Domitian, c.\ AD 95; an earlier date in the 60s is argued}
\dossiernote{The verse stands inside a rebuke: Laodicea is lukewarm, self-satisfied, and counselled to buy gold refined by fire (3:14--19). The knock is addressed to a community that thinks it needs nothing, and the promised meal is the offer that follows the reproof.}
\end{dossiertable}
